\chapter{Glucose Tolerance Test}

\section*{Overview}
A glucose tolerance test measures how the body processes a known glucose load by comparing blood glucose levels over time.

\section*{Common Instruments and Equipment}
The exact setup varies by technique and setting. Common items may include:
\begin{itemize}[nosep,leftmargin=*]
  \item glucose drink
  \item timed blood collection tubes and venipuncture supplies
  \item laboratory glucose analyzer
\end{itemize}

\section*{Why It Is Done}
It is used to evaluate diabetes or impaired glucose tolerance and is commonly used to screen for gestational diabetes during pregnancy.

\section*{How It Is Performed}
After fasting when required, a blood sample is collected, the patient drinks a measured glucose solution, and additional blood samples are taken at specified intervals.

\section*{Contraindications}
No single contraindication applies to every patient; the treating team assesses whether the procedure is suitable for the individual and the reason for treatment.

\section*{Precautions}
The care team reviews relevant health conditions, medicines, allergies, and procedure-specific risks, and explains any changes needed for the individual.

\section*{Preparation and Results}
Follow fasting, medication, and activity instructions. Nausea, dizziness, or weakness may occur after the glucose drink. Results are interpreted using the laboratory's criteria and the clinical context.

\section*{References}
\begin{enumerate}
  \item MedlinePlus. Glucose Tolerance Test. U.S. National Library of Medicine; 2025.
  \item Current Medical Diagnosis and Treatment. McGraw-Hill; 2025.
\end{enumerate}
\clearpage
